\section{Vanishing of the bracket tensors}\label{sec:vanishing}

We now test the bracket obstruction on sections supplied by deformations
of a rational curve. The tangent box makes the complementary target
bundle trivial along the curve where the test is performed. This use of
rational curves to force bracket vanishing has an antecedent in
\cite[Section~2.B and Lemma~3.18]{Hor07}.

Return to the hypotheses of \cref{thm:main}. For \(i\ne j\), define
\[
\BB_i:\bigwedge\nolimits^2E_i\longrightarrow E_j,\qquad
\BB_i(e,e')=P_j[e,e'].
\]
The extra terms in the bracket of scalar multiples lie in \(E_i\),
so \(P_j\) kills them. Thus \(\BB_i\) is a holomorphic bundle
morphism, and \(E_i\) is integrable exactly when \(\BB_i=0\).

The following consequence of rational connectedness gives the needed
sections. Its direct parameter-space proof is in
Appendix~\ref{app:variation}.

\begin{lemma}\label[lemma]{lem:variation}
Fix distinct \(a,b\in C\). There are a nonempty complex manifold
\(T\) and a holomorphic map \(G:C\times T\to X\) such that
\[
(\ev_a,\ev_b):T\longrightarrow X\times X,\qquad
\tau\longmapsto (G(a,\tau),G(b,\tau))
\]
is a submersion.
In particular, for \(g=G(\,\cdot\,,\tau)\), the global sections
of \(g^*E_i\) vanishing at \(a\) have values spanning \(E_i\) at
\(g(b)\).
\end{lemma}

\begin{proof}[Proof of \cref{thm:main}]
Choose the family in \cref{lem:variation}. For each member \(g\), take a
tangent box \(F:S\to X\) from \cref{prop:box}; the boxes need not be
chosen holomorphically as the member varies. For \(j\ne i\), consider
\[
L_i=\{(s_1,s_2)\in S:s_j=b\}.
\]
The free coordinate identifies \(L_i\) with \(C\). As shown in
\cref{fig:box}, its marked endpoint \((b,b)\) lies on the diagonal.
By \cref{prop:transport}, \(\Theta_i\) identifies
\(F^*E_i|_{L_i}\) with \(g^*E_i\), while \(\Theta_j\) identifies
\(F^*E_j|_{L_i}\) with the trivial bundle having fiber \(E_{j,g(b)}\).

Let \(u,v\) be any two global sections of \(g^*E_i\) vanishing at
\(a\). Restricting the transport maps to \(L_i\), the tensor value
\[
\sigma_{u,v}
=\Theta_j^{-1}\!\left((F^*\BB_i)(\Theta_i u\wedge\Theta_i v)\right)
\in H^0\bigl(C,\OO_C\otimes E_{j,g(b)}\bigr)
\]
is a holomorphic section of a trivial bundle on \(\PP^1\), hence is
constant. It vanishes at the point with free coordinate \(a\), so it
also vanishes at the point with free coordinate \(b\). At this latter
point, \(F(b,b)=g(b)\) and both transport maps are the identity.
Consequently
\[
\BB_i(g(b))\bigl(u(b)\wedge v(b)\bigr)=0.
\]
By \cref{lem:variation}, the values \(u(b)\) of sections vanishing at
\(a\) span \(E_{i,g(b)}\). Their pairwise wedges span its exterior
square, and therefore \(\BB_i(g(b))=0\).

This applies to both \(i\) and every member of the family. The map
\(\ev_b\) is a submersion, being the composition of the two-point
submersion with the second projection. Its image is a nonempty open
subset of \(X\), on which both bracket tensors vanish. The identity
principle for holomorphic bundle morphisms on connected \(X\) makes
them vanish everywhere. This is precisely integrability of both
specified summands.
\end{proof}

\begin{figure}[ht]
\centering
\begin{tikzpicture}[x=1.7cm,y=1.25cm,>=Stealth,
 every node/.style={font=\small}]
 \draw[gray!60] (0,0) rectangle (3,2.2);
 \draw[densely dashed,gray!75] (0,0)--(3,2.2)
 node[pos=.48,above left,text=black] {\(\Delta\)};
 \draw[very thick,blue!65!black,->] (0,2.2)--(2.8,2.2);
 \draw[very thick,blue!65!black,->] (3,0)--(3,2.05);
 \fill (0,0) circle (1.5pt);
 \fill (0,2.2) circle (1.5pt);
 \fill (3,0) circle (1.5pt);
 \fill (3,2.2) circle (1.8pt);
 \node[below left] at (0,0) {\((a,a)\)};
 \node[left] at (0,2.2) {\((a,b)\)};
 \node[below] at (3,0) {\((b,a)\)};
 \node[above right] at (3,2.2) {\((b,b)\)};
 \node[above] at (1.4,2.2) {\(L_1:\ s_2=b\)};
 \node[right] at (3,1.05) {\(L_2:\ s_1=b\)};
 \node at (1.15,.45) {\(F|_\Delta=g\)};
\end{tikzpicture}
\caption{The two factor curves used to test the bracket tensors,
drawn schematically in \(C\times C\).
On \(L_i\simeq\PP^1\), the input bundle is \(g^*E_i\) and the
complementary target bundle is trivial. Sections that vanish at
the marked off-diagonal point force their tensor values to
vanish at \((b,b)\), which maps to \(g(b)\).}
\label{fig:box}
\end{figure}


